\section*{Capítulo \ref{ch:b1:acetals-protection} --- Los carbonilos: adiciones nucleófilas, acetales y protección}

\begin{solution}{exo:b1:acetals-protection:1}
\ce{CH3CH(OH)OCH3}: \omterm{def:b1:acetals-protection:hemiacetal}{hemiacetal}. \ce{CH3CH(OCH3)2}: \omterm{def:b1:acetals-protection:hemiacetal}{acetal}.
\ce{CH3CH(OH)2}: hidrato. \ce{CH3OCH3}: éter. \ce{(CH3)2C(OCH2CH3)2}:
\omterm{def:b1:acetals-protection:hemiacetal}{acetal} (de una cetona).
\end{solution}

\begin{solution}{exo:b1:acetals-protection:2}
\ce{CH3CH2CHO + 2CH3OH <=> CH3CH2CH(OCH3)2 + H2O}; la propanona con
etano-1,2-diol da 2,2-dimetil-1,3-dioxolano y agua,
\ce{(CH3)2CO + HOCH2CH2OH <=> C5H10O2 + H2O}.
\end{solution}

\begin{solution}{exo:b1:acetals-protection:3}
El hidróxido de sodio, el bromuro de metilmagnesio y el borohidruro de
sodio lo dejan sin cambios; el ácido sulfúrico diluido en agua lo
hidroliza.
\end{solution}

\begin{solution}{exo:b1:acetals-protection:4}
El OH del carbono 4 se adiciona al carbono 1 del aldehído: un anillo de
cinco miembros, de cuatro carbonos y un oxígeno, con un OH en el carbono
contiguo al oxígeno del anillo.
\end{solution}

\begin{solution}{exo:b1:acetals-protection:5}
Protonación del C=O; el metanol se adiciona al carbono; pérdida de
\ce{H+}: \omterm{def:b1:acetals-protection:hemiacetal}{hemiacetal}; protonación de su OH; pérdida de agua hasta
\ce{CH3CH=O+CH3}; se adiciona un segundo metanol; pérdida de \ce{H+}:
\ce{CH3CH(OCH3)2}. Los dos protones tomados se devuelven: el ácido es un
\omterm{def:b1:elementary-steps:catalyst}{catalizador}.
\end{solution}

\begin{solution}{exo:b1:acetals-protection:6}
Protonación de uno de los OCH3; pérdida de metanol hasta
\ce{(CH3)2C=O+CH3}; se adiciona el agua; pérdida de \ce{H+}: \omterm{def:b1:acetals-protection:hemiacetal}{hemiacetal};
protonación, pérdida de metanol, pérdida de \ce{H+}: propanona. El exceso
de agua mantiene $Q$ por debajo de $K$ para la hidrólisis, que se completa.
\end{solution}

\begin{solution}{exo:b1:acetals-protection:7}
Una molécula de agua por \omterm{def:b1:acetals-protection:hemiacetal}{acetal}: $0.250 \times 18.0 = \qty{4.5}{g}$, unos
\qty{4.5}{mL}. Al cabo de una hora, $3.2/18.0 = \qty{0.178}{mol}$: una
conversión del \qty{71}{\percent}.
\end{solution}

\begin{solution}{exo:b1:acetals-protection:8}
Con el diol, una molécula sustituye a dos: la reacción no disminuye el
número de moléculas (un aldehído y un diol dan un \omterm{def:b1:acetals-protection:hemiacetal}{acetal} y un agua), y el
segundo OH se mantiene cerca del carbono que reacciona, lo que cierra el
anillo con facilidad. Ambas razones hacen más favorable el \omterm{def:b1:acetals-protection:hemiacetal}{acetal} cíclico.
\end{solution}

\begin{solution}{exo:b1:acetals-protection:9}
El ácido destruiría el \omterm{def:b1:organomagnesium:organometallic}{reactivo de Grignard} (transferencia de protón);
también hay que eliminar las trazas de agua y de diol.
\end{solution}

\begin{solution}{exo:b1:acetals-protection:10}
1. Se protege la cetona de la 4-clorobutan-2-ona como \omterm{def:b1:acetals-protection:hemiacetal}{acetal} cíclico
(etano-1,2-diol, \omterm{def:b1:elementary-steps:catalyst}{catalizador} ácido, Dean--Stark). 2. Magnesio en éter
seco: el \omterm{def:b1:organomagnesium:organometallic}{reactivo de Grignard} del cloruro protegido, que su propio
carbonilo ya no destruye. 3. Se añade propanona; tratamiento con cloruro
de amonio acuoso: el alcohol terciario. 4. Ácido acuoso diluido:
\omterm{def:b1:acetals-protection:protecting-group}{desprotección} a 5-hidroxi-5-metilhexan-2-ona.
\end{solution}

\begin{solution}{exo:b1:acetals-protection:11}
$Q = [\text{acetal}][\ce{H2O}]/([\text{aldehído}][\ce{CH3OH}]^2)$: un gran
exceso de metanol, elevado al cuadrado en el denominador, mantiene
$Q < K$ hasta que queda poco aldehído. Eliminar el agua actúa sobre el
numerador; usar el metanol como disolvente es más sencillo, pero exige un
producto fácil de separar.
\end{solution}

\begin{solution}{exo:b1:acetals-protection:12}
Solo el OH del \omterm{def:b1:acetals-protection:hemiacetal}{hemiacetal} puede salir como agua dando un catión
estabilizado por el oxígeno vecino del anillo (\ce{C=O+}); los demás
grupos OH tendrían que salir de carbonos ordinarios, dando cationes no
estabilizados. El metanol se adiciona al catión estabilizado: un \omterm{def:b1:acetals-protection:hemiacetal}{acetal}
de metilo.
\end{solution}

\begin{solution}{pb:b1:acetals-protection:1}
\textbf{1.} \ce{BrMgCH2CH2CH2CHO}.
\textbf{2.} Su carbono, de carácter carbaniónico, atacaría el aldehído de
otra molécula (o de sí misma): \ce{RMgBr} se adiciona a $\mathrm{R'CHO}$,
dando un \omterm{def:b1:alcohol-activation:alkoxide}{alcóxido}.
\textbf{3.} El aldehído, frente al \omterm{def:b1:organomagnesium:organometallic}{reactivo de Grignard}.
\textbf{4.} Los \omterm{def:b1:acetals-protection:hemiacetal}{acetales} son estables frente a los \omterm{def:b1:organomagnesium:organometallic}{reactivos de Grignard} y
a las bases, y se eliminan con ácido acuoso.
\textbf{5.} Se forma de manera más completa (un diol por dos moléculas de
alcohol) y es fácil de hidrolizar.
\textbf{6.} $15.09/150.9 = \qty{0.1000}{mol}$.
\textbf{7.} $1.20 \times 0.1000 \times 62.0 = \qty{7.44}{g}$.
\textbf{8.} \ce{BrCH2CH2CH2CHO + HOCH2CH2OH <=> C6H11BrO2 + H2O}
(2-(3-bromopropil)-1,3-dioxolano).
\textbf{9.} Activa el carbonilo y permite la pérdida de agua en la segunda
\omterm{def:b1:extent-q-and-k:system}{fase}; se regenera.
\textbf{10.} El tolueno y el agua destilan juntos, se condensan y caen en
el colector; el agua se hunde y se queda, y el tolueno rebosa de vuelta.
\textbf{11.} $Q$ contiene $[\ce{H2O}]$ en su numerador; eliminar el agua
mantiene $Q < K$, y la reacción avanza.
\textbf{12.} El agua es más densa y no se mezcla con el tolueno: se
acumula en el fondo; solo la capa superior llega al brazo de retorno.
\textbf{13.} Cuando el agua deja de acumularse, en el volumen esperado.
\textbf{14.} $0.1000 \times 194.9 = \qty{19.5}{g}$.
\textbf{15.} \ce{C6H11BrO2 + Mg -> C6H11BrMgO2} (éter seco).
\textbf{16.} El reactivo se adiciona a \ce{(CH3)2CO}: el \omterm{def:b1:alcohol-activation:alkoxide}{alcóxido}
\ce{(CH3)2C(O^-)CH2CH2CH2-} sobre la cadena protegida, con \ce{MgBr+}.
\textbf{17.} El ácido eliminaría el \omterm{def:b1:acetals-protection:protecting-group}{grupo protector} demasiado pronto y
podría deshidratar el alcohol terciario.
\textbf{18.} $0.0700 \times 174.0 = \qty{12.2}{g}$.
\textbf{19.} El \omterm{def:b1:acetals-protection:hemiacetal}{acetal} + \ce{H2O} $\to$ el aldehído + etano-1,2-diol
(\omterm{def:b1:elementary-steps:catalyst}{catalizador} ácido).
\textbf{20.} Su pérdida de agua exige un \omterm{def:b1:predominant-reaction:strong-acid}{ácido fuerte} y calor; el ácido
acuoso suave y frío hidroliza el \omterm{def:b1:acetals-protection:hemiacetal}{acetal} mucho más deprisa.
\textbf{21.} El OH del carbono 5 se adiciona al carbono 1 del aldehído: un
anillo de seis miembros (cinco carbonos y un oxígeno) con dos grupos
metilo en el carbono contiguo al oxígeno del anillo y un OH en el otro
carbono contiguo a él.
\textbf{22.} $0.0700 \times 0.90 = \qty{0.0630}{mol}$, \qty{8.19}{g}.
\textbf{23.} $0.0630/0.1000 = \qty{63}{\percent}$ (tomando la \omterm{def:b1:acetals-protection:protecting-group}{protección}
como completa).
\textbf{24.} $0.1000 \times 18.0 = \qty{1.80}{g}$.
\textbf{25.} $1.80/0.997 =$ \textbf{\qty{1.8}{mL} de agua} en el colector.
\end{solution}
